% 本书原创：按正文紧支集函数 beta 绘出三个频率副本；横坐标单位为 1/h。
\begin{figure}[htbp]
\centering
\begin{tikzpicture}[x=1.5cm,y=1cm,>=stealth,every node/.style={font=\small},line width=.65pt]
\begin{scope}
\fill[AnalysisBlue!8] (-.5,0) rectangle (.5,1.35);
\draw[gray,dashed] (-.5,0)--(-.5,1.35);
\draw[gray,dashed] (.5,0)--(.5,1.35);
\foreach \j in {-1,1}{
  \draw[orange!70!black,dashed,thick,domain=-.999:.999,samples=101,variable=\t]
    plot ({\j+.4*\t},{1.15*exp(1-1/(1-\t*\t))});
}
\draw[AnalysisBlue,thick,domain=-.999:.999,samples=101,variable=\t]
  plot ({.4*\t},{1.15*exp(1-1/(1-\t*\t))});
\draw[->,gray] (-1.84,0)--(1.84,0) node[above right] {\(\xi\)};
\foreach \j in {-1,0,1}{\draw[gray] (\j,.04)--(\j,-.04);}
\node[below] at (-1,-.05) {\(-1/h\)};
\node[below] at (0,-.05) {\(0\)};
\node[below] at (1,-.05) {\(1/h\)};
\node[below] at (-.5,-.56) {\(-1/(2h)\)};
\node[below] at (.5,-.56) {\(1/(2h)\)};
\draw[<->,gray] (0,1.55)--(1,1.55) node[midway,above] {\(1/h\)};
\draw[<->,gray] (-.5,-1.10)--(.5,-1.10) node[midway,below] {基本频带};
\node at (0,2.45) {副本分离：\(B=2/(5h)\)};
\end{scope}
\begin{scope}[xshift=6.5cm]
\fill[AnalysisBlue!8] (-.5,0) rectangle (.5,1.35);
\fill[orange!17] (-.5,0) rectangle (-.25,1.35);
\fill[orange!17] (.25,0) rectangle (.5,1.35);
\draw[gray,dashed] (-.5,0)--(-.5,1.35);
\draw[gray,dashed] (.5,0)--(.5,1.35);
\foreach \j in {-1,1}{
  \draw[orange!70!black,dashed,thick,domain=-.999:.999,samples=101,variable=\t]
    plot ({\j+.75*\t},{1.15*exp(1-1/(1-\t*\t))});
}
\draw[AnalysisBlue,thick,domain=-.999:.999,samples=101,variable=\t]
  plot ({.75*\t},{1.15*exp(1-1/(1-\t*\t))});
\draw[->,gray] (-1.84,0)--(1.84,0) node[above right] {\(\xi\)};
\foreach \j in {-1,0,1}{\draw[gray] (\j,.04)--(\j,-.04);}
\node[below] at (-1,-.05) {\(-1/h\)};
\node[below] at (0,-.05) {\(0\)};
\node[below] at (1,-.05) {\(1/h\)};
\node[below] at (-.5,-.56) {\(-1/(2h)\)};
\node[below] at (.5,-.56) {\(1/(2h)\)};
\draw[<->,gray] (0,1.55)--(1,1.55) node[midway,above] {\(1/h\)};
\draw[<->,gray] (-.5,-1.10)--(.5,-1.10) node[midway,below] {基本频带};
\node at (0,2.45) {副本重叠：\(B=3/(4h)\)};
\end{scope}
\node at (2.167,-1.90) {蓝实线：\(\Phi_B(\xi)\)；棕虚线：\(\Phi_B(\xi-1/h)\)、\(\Phi_B(\xi+1/h)\)};
\end{tikzpicture}
\caption[采样引起的频谱副本与混叠]{采样引起的频谱副本与混叠。按正文定义的 \(\Phi_B\) 绘出三个相邻副本，横轴采用相同的 \(1/h\) 单位，曲线采用同一纵向比例。竖虚线标出基本频带边界；右图橙色区域的内部是基本频带内两个副本同时非零的区间，曲线各自绘出，未画相加后的总和。}
\label{b2:fig:06-sampling-aliases}
\end{figure}
